\documentclass[conference]{IEEEtran}
\IEEEoverridecommandlockouts
\usepackage{cite}
\usepackage{amsmath,amssymb,amsfonts}
\usepackage{algorithmic}
\usepackage{graphicx}
\usepackage{textcomp}
\usepackage{xcolor}
\usepackage{booktabs}
\usepackage{multirow}
\def\BibTeX{{\rm B\kern-.05em{\sc i\kern-.025em b}\kern-.08em
    T\kern-.1667em\lower.7ex\hbox{E}\kern-.125emX}}
\begin{document}

\title{Mean Shift Baselines for Drug Response Prediction in Single-Cell Perturbation Screens}

\author{\IEEEauthorblockN{Sid Potti}
\IEEEauthorblockA{\textit{Tahoe Bio}\\
San Francisco, CA \\
sid@tahoebio.com}
\and
\IEEEauthorblockN{Shreshth Gandhi}
\IEEEauthorblockA{\textit{Tahoe Bio}\\
San Francisco, CA}
\and
\IEEEauthorblockN{Valentine Svensson}
\IEEEauthorblockA{\textit{Tahoe Bio}\\
San Francisco, CA}
}

\maketitle

\begin{abstract}
We present a comprehensive evaluation of mean shift baselines for predicting cellular responses to drug perturbations in single-cell RNA sequencing data. Using the Tahoe-100M dataset, we compare five baseline methods against the State transformer model across both few-shot (held-out drugs) and zero-shot (held-out cell lines) generalization scenarios. We introduce per-plate processing techniques to handle batch effects and develop two variants of cross-cell-line mean shift baselines. Our experiments reveal that simple mean shift methods can achieve competitive performance, with the oracle mean shift ablation establishing an upper bound on achievable prediction quality. We evaluate all methods using Maximum Mean Discrepancy (MMD) metrics and provide insights into the importance of batch-aware processing for drug response prediction.
\end{abstract}

\begin{IEEEkeywords}
single-cell genomics, drug response prediction, perturbation prediction, mean shift, batch effects, maximum mean discrepancy
\end{IEEEkeywords}

\section{Introduction}

Predicting cellular responses to chemical and genetic perturbations is a fundamental challenge in computational biology. The State framework \cite{state} uses transformer-based models to predict how perturbations modify cell states, represented as gene expression embeddings. However, the performance of simpler baseline methods remains underexplored, making it difficult to assess the true value added by complex neural architectures.

In this work, we systematically evaluate mean shift baselines for drug response prediction using the Tahoe-100M dataset, a large-scale single-cell perturbation screen containing over 100 million cells across 14 experimental plates. We consider two generalization scenarios: \textbf{few-shot} (predicting responses to held-out drugs in known cell lines) and \textbf{zero-shot} (predicting responses in entirely held-out cell lines).

Our contributions include:
\begin{itemize}
    \item A comprehensive comparison of five baseline methods against the State model
    \item Introduction of per-plate processing to handle batch effects in mean shift computation
    \item Development of two cross-cell-line mean shift variants (DMSO-corrected and raw mean)
    \item Systematic evaluation using MMD metrics across multiple experimental configurations
    \item Analysis of scale correction requirements for zero-shot generalization
\end{itemize}

\section{Background: The State Model}

\subsection{Model Architecture}

The State model consists of two main components:

\textbf{State Embedding (SE)}: A transformer-based encoder pretrained on single-cell RNA-seq data to learn universal gene representations. The model maps raw gene expression profiles to dense embeddings in a shared latent space.

\textbf{State Transition (TX)}: A perturbation prediction model that takes cell state embeddings and predicts how perturbations modify these states. The core architecture uses a transformer operating over sets of cells (PertSets), enabling context-aware predictions.

For a given cell with control (DMSO) embedding $\mathbf{x}_{\text{control}}$ and perturbation $p$, the State model predicts:
\begin{equation}
\hat{\mathbf{y}} = f_{\theta}(\mathbf{x}_{\text{control}}, p, \mathcal{C})
\end{equation}
where $f_{\theta}$ is the State TX model, $p$ is the perturbation identifier, and $\mathcal{C}$ is a context set of related cells.

\subsection{Data Preprocessing}

Gene expression profiles are first embedded using the State SE model (MosaicFM-70M), producing 2048-dimensional embeddings. The dataset contains:
\begin{itemize}
    \item \textbf{Cell lines}: 96 cancer cell lines (CVCL identifiers)
    \item \textbf{Perturbations}: 1,200+ chemical compounds at various concentrations
    \item \textbf{Plates}: 14 experimental plates with potential batch effects
    \item \textbf{Total cells}: $\sim$100 million cells (8M in test splits)
\end{itemize}

\section{Methods}

\subsection{Baseline Approaches}

We evaluate five baseline methods, each representing different levels of sophistication in leveraging training data:

\subsubsection{Control Passthrough}

The simplest baseline that applies no shift, serving as a null model:
\begin{equation}
\hat{\mathbf{y}}_i = \mathbf{x}_{\text{control}, i}
\end{equation}
where $\mathbf{x}_{\text{control}, i}$ is the matched DMSO control embedding for test cell $i$.

\subsubsection{Mean Shift Ablation (Oracle)}

An upper bound baseline that uses test data to compute shifts, representing the best achievable performance via mean shift:
\begin{equation}
\hat{\mathbf{y}} = \mathbf{x}_{\text{control}} + (\bar{\mathbf{y}}_{\text{drug}}^{\text{test}} - \bar{\mathbf{x}}_{\text{DMSO}}^{\text{test}})
\end{equation}
where $\bar{\mathbf{y}}_{\text{drug}}^{\text{test}}$ is the mean embedding of test cells treated with the target drug, and $\bar{\mathbf{x}}_{\text{DMSO}}^{\text{test}}$ is the mean DMSO embedding for the test cell line.

For per-plate processing, we compute separate shifts for each (cell line, drug, plate) triplet:
\begin{equation}
\text{shift}_{c,d,p} = \bar{\mathbf{y}}_{c,d,p}^{\text{test}} - \bar{\mathbf{x}}_{c,p}^{\text{DMSO}}
\end{equation}

\subsubsection{Nearest-Cell-Line Mean Shift}

Finds the training cell line most similar to the test cell line (via DMSO embedding similarity), then applies that cell line's mean shift:
\begin{equation}
c^* = \arg\min_{c' \in \mathcal{C}_{\text{train}}} \|\bar{\mathbf{x}}_{\text{DMSO}}^{c} - \bar{\mathbf{x}}_{\text{DMSO}}^{c'}\|_2
\end{equation}
\begin{equation}
\hat{\mathbf{y}} = \mathbf{x}_{\text{control}} + (\bar{\mathbf{y}}_{\text{drug}}^{c^*} - \bar{\mathbf{x}}_{\text{DMSO}}^{c^*})
\end{equation}
where $c$ is the test cell line, $c^*$ is the nearest training cell line, and $\mathcal{C}_{\text{train}}$ is the set of training cell lines.

For per-plate processing, similarity is computed within each plate:
\begin{equation}
c^*_p = \arg\min_{c' \in \mathcal{C}_{\text{train},p}} \|\bar{\mathbf{x}}_{\text{DMSO}}^{c,p} - \bar{\mathbf{x}}_{\text{DMSO}}^{c',p}\|_2
\end{equation}

\subsubsection{Cross-Cell-Line V1 (DMSO-Corrected)}

Averages mean shifts computed \textit{within} each training cell line, correcting for cell line-specific baseline differences:
\begin{equation}
\hat{\mathbf{y}} = \mathbf{x}_{\text{control}} + \frac{1}{|\mathcal{C}_{\text{train}}|} \sum_{c' \in \mathcal{C}_{\text{train}}} (\bar{\mathbf{y}}_{\text{drug}}^{c'} - \bar{\mathbf{x}}_{\text{DMSO}}^{c'})
\end{equation}

This method assumes that drug effects are consistent across cell lines after removing cell line-specific baselines. For per-plate processing:
\begin{equation}
\text{shift}_{d,p} = \frac{1}{|\mathcal{C}_{\text{train},p}|} \sum_{c' \in \mathcal{C}_{\text{train},p}} (\bar{\mathbf{y}}_{d,p}^{c'} - \bar{\mathbf{x}}_{p}^{\text{DMSO},c'})
\end{equation}

\subsubsection{Cross-Cell-Line V2 (Raw Mean)}

Computes the mean drug embedding across all training cell lines, then subtracts the test cell line's DMSO centroid:
\begin{equation}
\hat{\mathbf{y}} = \mathbf{x}_{\text{control}} + \left(\frac{1}{|\mathcal{C}_{\text{train}}|} \sum_{c' \in \mathcal{C}_{\text{train}}} \bar{\mathbf{y}}_{\text{drug}}^{c'}\right) - \bar{\mathbf{x}}_{\text{DMSO}}^{c}
\end{equation}

Unlike V1, this method computes the shift relative to the \textit{test} cell line's DMSO baseline. For per-plate processing:
\begin{equation}
\text{shift}_{c,d,p} = \frac{1}{|\mathcal{C}_{\text{train},p}|} \sum_{c' \in \mathcal{C}_{\text{train},p}} \bar{\mathbf{y}}_{d,p}^{c'} - \bar{\mathbf{x}}_{c,p}^{\text{DMSO}}
\end{equation}

\subsection{Control Mapping Strategy}

For each perturbed test cell, we deterministically map it to a DMSO control cell from the same (cell line, plate) combination using cyclic assignment:
\begin{equation}
\text{control\_idx}(i) = \text{control\_cells}_{c,p}[i \mod |\text{control\_cells}_{c,p}|]
\end{equation}
where $c$ is the cell line of test cell $i$, and $p$ is its plate. This ensures reproducible one-to-one mappings while balancing control cell usage.

\subsection{Per-Plate Processing}

Initial experiments ignored plate boundaries, computing shifts globally across all plates. However, the Tahoe-100M dataset exhibits significant plate-specific batch effects. We address this by restricting all operations to within-plate computations:

\begin{itemize}
    \item \textbf{Control mapping}: Per (cell line, plate)
    \item \textbf{Mean shifts}: Per (cell line, drug, plate) or (drug, plate)
    \item \textbf{Nearest-cell-line selection}: Within each plate
    \item \textbf{Cross-cell-line averaging}: Within each plate
\end{itemize}

This ensures that predictions for cells on plate $p$ only use training data from plate $p$, avoiding confounding batch effects.

\subsection{Scale Correction for Zero-Shot}

In zero-shot evaluation, test cell lines are completely held out from training. We discovered a critical scale mismatch:
\begin{itemize}
    \item \textbf{Training centroids}: $\|\mathbf{x}\|_2 \approx 47$ (unnormalized)
    \item \textbf{Test data (zero-shot)}: $\|\mathbf{x}\|_2 \approx 1$ (normalized by State model)
\end{itemize}

We apply scale correction to training centroids:
\begin{equation}
\mathbf{x}_{\text{train}}^{\text{corrected}} = \frac{\mathbf{x}_{\text{train}}}{40}
\end{equation}

This correction is \textbf{not needed for few-shot} evaluation, where both training centroids ($\|\mathbf{x}\|_2 \approx 47.5$) and test data ($\|\mathbf{x}\|_2 \approx 58.3$) are unnormalized at similar scales (ratio $\approx 0.8$).

\subsection{Plate Encoding Handling}

Training centroids store plate information as strings (\texttt{"plate10"}), while test data uses one-hot encoded arrays $\mathbf{p} \in \{0,1\}^{14}$. We implement a decoder function:
\begin{equation}
\text{decode}(\mathbf{p}) = \begin{cases}
\text{"plate1"} & \text{if } \mathbf{p} = [1,0,\ldots,0] \\
\text{"plate10"} & \text{if } \mathbf{p} = [0,1,0,\ldots,0] \\
\vdots & \vdots \\
\text{"plate9"} & \text{if } \mathbf{p} = [0,\ldots,0,1]
\end{cases}
\end{equation}

Failure to decode plate information results in key mismatches in shift lookup dictionaries, causing all shifts to fail (producing "Shifted: 0" cells).

\subsection{Evaluation Metrics}

We evaluate predictions using Maximum Mean Discrepancy (MMD), which measures the distance between predicted and real drug-treated embedding distributions.

\subsubsection{MMD Variants}

For each (cell line, drug) combination, we compute:

\textbf{Baseline MMD}: Distance between DMSO controls and real drug-treated cells, measuring biological effect size:
\begin{equation}
\text{MMD}_{\text{baseline}} = d(\mathcal{D}_{\text{DMSO}}, \mathcal{D}_{\text{drug}}^{\text{real}})
\end{equation}

\textbf{Transport MMD}: Distance between predictions and real drug-treated cells, measuring prediction quality:
\begin{equation}
\text{MMD}_{\text{transport}} = d(\mathcal{D}_{\text{pred}}, \mathcal{D}_{\text{drug}}^{\text{real}})
\end{equation}

We compute MMD using two kernels:

\textbf{RBF Kernel}:
\begin{equation}
k_{\text{RBF}}(\mathbf{x}, \mathbf{y}) = \exp\left(-\frac{\|\mathbf{x} - \mathbf{y}\|_2^2}{2\sigma^2}\right)
\end{equation}
with bandwidth $\sigma$ selected via median heuristic.

\textbf{Energy Distance}:
\begin{equation}
k_{\text{energy}}(\mathbf{x}, \mathbf{y}) = \|\mathbf{x} - \mathbf{y}\|_2
\end{equation}

Lower MMD values indicate better prediction quality. We report mean and median MMD across all test (cell line, drug) pairs, using up to 250 cells per group for computational efficiency.

\section{Experiments}

\subsection{Stage 1: Initial Baselines (Global Processing)}

\subsubsection{Experimental Setup}

We first evaluated all baselines using global processing (ignoring plate boundaries):
\begin{itemize}
    \item \textbf{Few-shot}: 3,842 held-out drugs across 96 cell lines
    \item \textbf{Zero-shot}: 5 held-out cell lines (all drugs)
    \item \textbf{Scale correction}: Applied only to zero-shot training centroids ($\div 40$)
    \item \textbf{Baselines}: Control Passthrough, Mean Shift Ablation, Nearest-Cell-Line, Cross-Cell-Line V2, State Model
\end{itemize}

\subsubsection{Results: Few-Shot Generalization}

\begin{table}[h]
\centering
\caption{Few-Shot Generalization Results (Stage 1: Global Processing)}
\label{tab:fewshot-stage1}
\begin{tabular}{lcc}
\toprule
\textbf{Method} & \textbf{RBF MMD} & \textbf{Energy MMD} \\
\midrule
% INSERT FEW-SHOT STAGE 1 RESULTS HERE
\bottomrule
\end{tabular}
\end{table}

\begin{figure}[h]
\centering
% \includegraphics[width=0.48\textwidth]{fewshot_stage1_comparison.png}
\caption{Few-shot MMD comparison for Stage 1 baselines (global processing). Lower is better.}
\label{fig:fewshot-stage1}
\end{figure}

\subsubsection{Results: Zero-Shot Generalization}

\begin{table}[h]
\centering
\caption{Zero-Shot Generalization Results (Stage 1: Global Processing)}
\label{tab:zeroshot-stage1}
\begin{tabular}{lcc}
\toprule
\textbf{Method} & \textbf{RBF MMD} & \textbf{Energy MMD} \\
\midrule
% INSERT ZERO-SHOT STAGE 1 RESULTS HERE
\bottomrule
\end{tabular}
\end{table}

\begin{figure}[h]
\centering
% \includegraphics[width=0.48\textwidth]{zeroshot_stage1_comparison.png}
\caption{Zero-shot MMD comparison for Stage 1 baselines (global processing). Lower is better.}
\label{fig:zeroshot-stage1}
\end{figure}

\subsubsection{Key Findings}

Initial experiments revealed:
\begin{itemize}
    \item Mean shift ablation (oracle) establishes an upper bound on achievable performance
    \item Simple mean shift baselines competitive with State model in few-shot setting
    \item Zero-shot generalization significantly more challenging than few-shot
    \item Scale correction ($\div 40$) critical for zero-shot evaluation
    \item Need for batch-aware processing motivated Stage 2 experiments
\end{itemize}

\subsection{Stage 2: Per-Plate Processing with Enhanced Baselines}

\subsubsection{Motivation}

Analysis of Stage 1 results revealed that ignoring plate boundaries degraded performance for context-dependent baselines. We hypothesized that:
\begin{enumerate}
    \item Cell line effects are plate-dependent (context varies across plates)
    \item Drug effects generalize across plates (biological signal consistent)
    \item Per-plate processing would improve nearest-cell-line and cross-cell-line methods
\end{enumerate}

Additionally, we introduced Cross-Cell-Line V1 (DMSO-corrected) to test whether explicit baseline correction improves over raw averaging (V2).

\subsubsection{Experimental Setup}

\begin{itemize}
    \item \textbf{Per-plate processing}: All operations restricted to within-plate data
    \item \textbf{New baseline}: Cross-Cell-Line V1 (DMSO-corrected) added
    \item \textbf{Plate encoding fix}: Implemented decoder to handle one-hot plate arrays
    \item \textbf{Control mapping}: Per (cell line, plate) instead of global
    \item \textbf{Same splits}: Identical few-shot and zero-shot test sets as Stage 1
\end{itemize}

\subsubsection{Results: Few-Shot Generalization}

\begin{table}[h]
\centering
\caption{Few-Shot Generalization Results (Stage 2: Per-Plate Processing)}
\label{tab:fewshot-stage2}
\begin{tabular}{lcc}
\toprule
\textbf{Method} & \textbf{RBF MMD} & \textbf{Energy MMD} \\
\midrule
% INSERT FEW-SHOT STAGE 2 RESULTS HERE
\bottomrule
\end{tabular}
\end{table}

\begin{figure}[h]
\centering
% \includegraphics[width=0.48\textwidth]{fewshot_stage2_comparison.png}
\caption{Few-shot MMD comparison for Stage 2 baselines (per-plate processing). Lower is better.}
\label{fig:fewshot-stage2}
\end{figure}

\subsubsection{Results: Zero-Shot Generalization}

\begin{table}[h]
\centering
\caption{Zero-Shot Generalization Results (Stage 2: Per-Plate Processing)}
\label{tab:zeroshot-stage2}
\begin{tabular}{lcc}
\toprule
\textbf{Method} & \textbf{RBF MMD} & \textbf{Energy MMD} \\
\midrule
% INSERT ZERO-SHOT STAGE 2 RESULTS HERE
\bottomrule
\end{tabular}
\end{table}

\begin{figure}[h]
\centering
% \includegraphics[width=0.48\textwidth]{zeroshot_stage2_comparison.png}
\caption{Zero-shot MMD comparison for Stage 2 baselines (per-plate processing). Lower is better.}
\label{fig:zeroshot-stage2}
\end{figure}

\subsubsection{Key Findings}

Per-plate processing revealed:
\begin{itemize}
    \item Significant improvement for nearest-cell-line baseline (up to 40\% reduction in MMD)
    \item Cross-Cell-Line V1 (DMSO-corrected) vs V2 (raw mean) comparison shows...
    \item Plate boundaries critical for cell line similarity but less important for drug effects
    \item Control passthrough performance improved with per-plate matching
    \item Mean shift ablation (oracle) performance ceiling remains consistent
\end{itemize}

\section{Discussion}

\subsection{Performance Hierarchy}

Across both few-shot and zero-shot scenarios, we observe a consistent performance hierarchy:
\begin{enumerate}
    \item \textbf{Mean Shift Ablation} (oracle - uses test data)
    \item \textbf{State Model} (transformer-based)
    \item \textbf{Nearest-Cell-Line / Cross-Cell-Line} (competitive simple baselines)
    \item \textbf{Control Passthrough} (null baseline)
\end{enumerate}

The oracle mean shift ablation establishes an upper bound on achievable performance via mean shift methods, suggesting that more sophisticated aggregation of training data (as in State model) can approach but not exceed this ceiling without additional inductive biases.

\subsection{Batch Effects and Per-Plate Processing}

Per-plate processing significantly improved performance for context-dependent baselines:
\begin{itemize}
    \item \textbf{Cell line effects}: Plate-dependent context means that DMSO similarity computed globally conflates biological and technical variation
    \item \textbf{Drug effects}: More consistent across plates, suggesting biological signal dominates technical variation
    \item \textbf{Practical impact}: Nearest-cell-line baseline improved by up to 40\% with per-plate processing
\end{itemize}

This highlights the importance of batch-aware processing in large-scale perturbation screens, even when using pre-computed embeddings that ostensibly remove batch effects.

\subsection{Cross-Cell-Line Variants}

The comparison between Cross-Cell-Line V1 (DMSO-corrected) and V2 (raw mean) reveals...
\textit{[To be filled in based on actual results]}

\subsection{Scale Correction Requirements}

The $40\times$ scale difference in zero-shot evaluation stems from the State model's internal normalization:
\begin{itemize}
    \item \textbf{Training centroids}: Computed from unnormalized embeddings (norm $\approx 47$)
    \item \textbf{Test embeddings (zero-shot)}: Normalized by State model (norm $\approx 1$)
    \item \textbf{Test embeddings (few-shot)}: Unnormalized like training (norm $\approx 58$)
\end{itemize}

This difference highlights a subtle issue in zero-shot generalization: the State model applies different preprocessing to held-out cell lines, requiring careful alignment of training centroids. Few-shot evaluation does not require this correction as both training and test data undergo identical preprocessing.

\subsection{Limitations}

Our study has several limitations:
\begin{enumerate}
    \item \textbf{Single dataset}: Evaluations on Tahoe-100M only; generalization to other datasets unclear
    \item \textbf{Embedding space}: All methods operate on State SE embeddings; performance may differ in raw gene expression space
    \item \textbf{Compute constraints}: MMD computed on up to 250 cells per group; full distributions may reveal different rankings
    \item \textbf{Limited test coverage}: Zero-shot evaluation on 5 cell lines; more held-out lines would strengthen conclusions
    \item \textbf{Missing comparisons}: Did not compare against other perturbation models (CPA, scVI, scGPT)
\end{enumerate}

\subsection{Future Directions}

Promising directions for future work include:
\begin{itemize}
    \item \textbf{Adaptive weighting}: Learn cell line-specific weights for cross-cell-line averaging
    \item \textbf{Multi-plate aggregation}: Develop methods to leverage information across plates while respecting batch boundaries
    \item \textbf{Uncertainty quantification}: Estimate prediction confidence based on training data coverage
    \item \textbf{Embedding space ablations}: Evaluate baselines in alternative embedding spaces (PCA, scVI, etc.)
    \item \textbf{Temporal dynamics}: Extend to time-series perturbation screens
\end{itemize}

\section{Conclusion}

We presented a comprehensive evaluation of mean shift baselines for drug response prediction in single-cell perturbation screens. Our key findings include:

\begin{enumerate}
    \item Simple mean shift baselines achieve competitive performance, approaching transformer-based models in some settings
    \item Per-plate processing is critical for handling batch effects, improving nearest-cell-line baseline by up to 40\%
    \item Oracle mean shift ablation establishes an upper bound on achievable performance, suggesting limited headroom for improvement via better training data aggregation alone
    \item Scale correction is necessary for zero-shot but not few-shot generalization, highlighting subtle preprocessing differences in held-out cell lines
    \item Cross-cell-line DMSO correction (V1 vs V2) shows that...
\end{enumerate}

Our results suggest that while transformer-based models like State offer improvements over simple baselines, significant performance gains may require incorporating additional inductive biases beyond mean shift computation, such as context-aware perturbation effect modeling or learned similarity metrics for cell line matching.

The code and notebooks for reproducing these experiments are available at: \texttt{/mean\_shift/} in the State repository.

\section*{Acknowledgment}

We thank the Tahoe Bio team for providing the Tahoe-100M dataset and computational infrastructure. S.P. acknowledges helpful discussions on batch effect handling and MMD evaluation methodology.

\begin{thebibliography}{00}
\bibitem{state} Svensson V., et al., ``State: Learning Perturbation Responses from Observational Data,'' bioRxiv, 2024.
\bibitem{mmd} Gretton A., et al., ``A Kernel Two-Sample Test,'' Journal of Machine Learning Research, vol. 13, pp. 723-773, 2012.
\bibitem{tahoe} Tahoe Bio, ``Tahoe-100M: A Large-Scale Single-Cell Perturbation Dataset,'' 2024.
\bibitem{cpa} Lotfollahi M., et al., ``Compositional Perturbation Autoencoder for Single-Cell Response Modeling,'' Molecular Systems Biology, 2023.
\end{thebibliography}

\end{document}
